\documentclass[12pt,a4paper]{article}

% ── Packages ──────────────────────────────────────────────────────────────────
\usepackage[a4paper, margin=2.5cm]{geometry}
\usepackage[T1]{fontenc}
\usepackage[utf8]{inputenc}
\usepackage{lmodern}
\usepackage{amsmath, amssymb}
\usepackage{booktabs}
\usepackage{array}
\usepackage{tabularx}
\usepackage{longtable}
\usepackage{parskip}
\usepackage{hyperref}
\usepackage{xcolor}
\usepackage{titlesec}
\usepackage{enumitem}
\usepackage{listings}
\usepackage{microtype}
\usepackage{fancyhdr}
\usepackage{lastpage}
\usepackage{graphicx}
\usepackage{float}

% ── Hyperref setup ────────────────────────────────────────────────────────────
\hypersetup{
    colorlinks=true,
    linkcolor=blue!70!black,
    urlcolor=blue!70!black,
    citecolor=blue!70!black,
    pdftitle={COS 710 Assignment 2 -- Structure-Based Genetic Programming for Regression},
    pdfauthor={Yohali Malaika Kamangu}
}

% ── Header / Footer ───────────────────────────────────────────────────────────
\pagestyle{fancy}
\fancyhf{}
\rhead{COS 710 -- Assignment 2}
\lhead{Yohali Malaika Kamangu (u23618583)}
\cfoot{\thepage\ of \pageref{LastPage}}
\renewcommand{\headrulewidth}{0.4pt}

% ── Section title formatting ──────────────────────────────────────────────────
\titleformat{\section}{\large\bfseries}{}{0em}{\thesection.\quad}
\titleformat{\subsection}{\normalsize\bfseries}{}{0em}{\thesubsection\quad}

% ── Code listing style ────────────────────────────────────────────────────────
\lstset{
    basicstyle=\ttfamily\small,
    breaklines=true,
    frame=single,
    backgroundcolor=\color{gray!8},
    rulecolor=\color{gray!40},
    xleftmargin=1em,
    xrightmargin=1em
}

% ── Column type for tables ────────────────────────────────────────────────────
\newcolumntype{L}[1]{>{\raggedright\arraybackslash}p{#1}}
\newcolumntype{C}[1]{>{\centering\arraybackslash}p{#1}}

% ──────────────────────────────────────────────────────────────────────────────
\begin{document}

% ── Title Page ────────────────────────────────────────────────────────────────
\begin{titlepage}
    \centering
    \vspace*{2cm}
    {\LARGE\bfseries COS 710: Artificial Intelligence\par}
    \vspace{0.5cm}
    {\Large Assignment 2: Structure-Based Genetic Programming for Regression\par}
    \vspace{1.5cm}
    {\large Yohali Malaika Kamangu\par}
    {\large u23618583\par}
    \vspace{0.5cm}
    {\large April 2026\par}
    \vspace{1.5cm}
    {\normalsize Dataset: Residential Energy Usage UK (2014--2020)\par}
\end{titlepage}

\tableofcontents
\newpage

% ──────────────────────────────────────────────────────────────────────────────
\section{Description of the Dataset}

The dataset originates from the UK residential energy sector and covers a nearly
five-year period from 31 December 2014 to 2020. It contains 201,604 records
collected at 15-minute intervals. The following attributes are recorded for each
timestamp:

\begin{center}
\begin{tabular}{@{}L{6.5cm}L{5.5cm}C{1.5cm}@{}}
\toprule
\textbf{Column} & \textbf{Description} & \textbf{Unit} \\
\midrule
\texttt{utc\_timestamp} & Date and time of the record (UTC) & -- \\
\texttt{Electricity\_load} & Residential electricity demand & kW \\
\texttt{Residential\_electricity\_price} & Residential electricity price & GBP \\
\texttt{Residential\_solar\_generation} & Solar power generation & kW \\
\texttt{Residential\_wind\_generation} & Wind power generation & kW \\
\texttt{Temperature} & Ambient temperature & °C \\
\texttt{Relative Humidity} & Relative humidity & \% \\
\bottomrule
\end{tabular}
\end{center}

Only the \texttt{Electricity\_load} column is used for regression. The problem is
formulated as univariate time-series prediction: given the electricity load values
at the $n$ most recent 15-minute intervals, predict the load at the next interval.
A sliding window of size 7 (105 minutes of history) is applied:

\begin{itemize}
  \item Input variables: $x_1 = \text{Load}(t{-}1),\; x_2 = \text{Load}(t{-}2),\; \ldots,\; x_7 = \text{Load}(t{-}7)$
  \item Target: $\text{Load}(t)$
\end{itemize}

To keep runtimes within a practical range while retaining sufficient temporal
coverage, the first 100,000 valid rows are loaded. This sub-sample spans
approximately 2.5 years of the 5-year dataset (2014--mid 2017), preserving the
circadian and seasonal patterns necessary for autoregressive prediction. The
resulting 99,993 windowed examples provide ample data for GP training: an 80/20
temporal split produces 79,994 training examples and 19,999 test examples. No
shuffling is applied, as shuffling would introduce data leakage in a time-series
setting.

Electricity load values lie predominantly in $[0.04,\, 0.15]$ kW with occasional
higher peak readings. The series exhibits clear circadian and weekly seasonality
typical of residential consumption patterns.

% ──────────────────────────────────────────────────────────────────────────────
\section{Description of the Performance Metric}

\textbf{Mean Squared Error (MSE)} is the primary fitness and evaluation metric:

\[
\text{MSE} = \frac{1}{n} \sum_{i=1}^{n} (\hat{y}_i - y_i)^2
\]

where $\hat{y}_i$ is the tree's prediction for example $i$ and $y_i$ is the true
load value. MSE is used for training fitness, test evaluation, and summary statistics
because it penalises large errors superlinearly and is consistent across runs.
$\text{RMSE} = \sqrt{\text{MSE}}$ is noted in discussion to interpret results in
kilowatts (kW).

\textbf{Parsimony pressure} is incorporated into the training fitness to discourage
bloated trees:

\[
\text{rawFitness}(T) = \text{MSE}(T) \times \bigl(1 + \lambda \cdot |T|\bigr), \quad \lambda = 0.001
\]

where $|T|$ denotes the node count of tree $T$.

\textbf{Hits criterion:} A prediction is a ``hit'' if $|\hat{y}_i - y_i| < 0.01$~kW
(approximately 10\% of a typical reading). The hit count for the best individual
is recorded per generation.

\textbf{Success predicate:} A run is successful if the best training MSE
falls below 0.00020.

% ──────────────────────────────────────────────────────────────────────────────
\section{How Structure is Incorporated into Canonical GP}

\subsection{Motivation}

Canonical GP directs the search solely through behavioural fitness -- how well a
tree's outputs match the training targets. This tends to produce structurally
homogeneous populations: once a particular operator--terminal pattern proves
beneficial it proliferates rapidly, reducing the diversity needed for continued
exploration. Assignment 2 addresses this by augmenting the fitness function with
a structural diversity component, following the first option described in the SBGP
lecture slides: \emph{``As part of the fitness.'''}

\subsection{The Similarity Index}

Following the lecture slide example, structural similarity is measured by
\textbf{counting the number of nodes that are the same from the root} -- a
top-down prefix node matching scheme:

\begin{lstlisting}[language=Java, caption={Prefix match count (Node.java)}]
public static int prefixMatchCount(Node n1, Node n2) {
    if (n1 == null || n2 == null) return 0;
    if (!n1.label.equals(n2.label)) return 0;
    int count = 1;
    int numChildren = Math.min(n1.children.size(), n2.children.size());
    for (int i = 0; i < numChildren; i++)
        count += prefixMatchCount(n1.children.get(i), n2.children.get(i));
    return count;
}
\end{lstlisting}

This count is normalised by the size of the larger tree to yield a dimensionless
similarity score in $[0,\, 1]$:

\[
\text{sim}(T_1, T_2) \;=\; \frac{\text{prefixMatchCount}(\text{root}_1,\, \text{root}_2)}{\max(|T_1|,\, |T_2|)}
\]

\begin{itemize}
  \item Two identical trees return 1.0.
  \item Two trees with different root labels return 0.0.
  \item Trees \texttt{(+ x1 x2)} and \texttt{(+ x1 x3)} return $2/3 \approx 0.667$
        (root \texttt{+} and left child \texttt{x1} match; right children differ).
\end{itemize}

\subsection{Structural Fitness Sharing}

After computing standard MSE-based raw fitness each generation, each individual's
fitness is adjusted upward (worsened) in proportion to how many structurally
similar peers it has in the population:

\[
\text{adjustedFitness}[i] \;=\; \text{rawFitness}[i] \times \bigl(1 + \alpha \cdot \text{nicheCount}[i]\bigr)
\]

where:
\begin{itemize}
  \item $\text{nicheCount}[i]$ = number of sampled peers $j$ with $\text{sim}(T_i, T_j) \geq \sigma$
  \item $\sigma = 0.60$ is the \textbf{similarity threshold}: a tree must share at
        least 60\% of its top-down node structure with individual $i$ to be counted
        as a structural neighbour.
  \item $\alpha = 0.01$ is the \textbf{sharing coefficient}: scales the penalty so
        the structural component complements rather than dominates the behavioural
        fitness signal.
\end{itemize}

To avoid the $O(n^2)$ cost of all-pairs comparison, each individual's niche count
is estimated by sampling 50 random peers, giving 50,000 structural comparisons
per generation at a fraction of the full computation cost.

\subsection{Selection Using Adjusted Fitness}

Tournament selection and elitism both use the \emph{adjusted} fitness. The
Fisher-Yates slot array determines \emph{which} operator (crossover or mutation)
is applied to each non-elite slot; the parents for those slots are then selected
via tournament on \emph{adjusted} fitness. The best individual is tracked and
reported by \emph{raw} fitness, since raw fitness (MSE + parsimony) reflects
actual prediction quality without structural inflation.

% ──────────────────────────────────────────────────────────────────────────────
\section{Representation, Terminal Set, and Function Set}

\subsection{Representation}

Each GP individual is a \textbf{full expression tree} rooted at a binary function
node. Nodes store a label (operator or terminal), arity, an ordered children list,
and a parent pointer.

\subsection{Terminal Set}

\begin{center}
\begin{tabular}{@{}ll@{}}
\toprule
\textbf{Terminal} & \textbf{Meaning} \\
\midrule
$x_1$ & $\text{Load}(t{-}1)$ -- most recent lag \\
$x_2$ & $\text{Load}(t{-}2)$ \\
$x_3$ & $\text{Load}(t{-}3)$ \\
$x_4$ & $\text{Load}(t{-}4)$ \\
$x_5$ & $\text{Load}(t{-}5)$ \\
$x_6$ & $\text{Load}(t{-}6)$ \\
$x_7$ & $\text{Load}(t{-}7)$ -- oldest lag \\
ERC   & Ephemeral Random Constant: $U(-5.0,\,5.0)$, 30\% prob.\ at terminal creation \\
\bottomrule
\end{tabular}
\end{center}

\subsection{Function Set}

\begin{center}
\begin{tabular}{@{}lll@{}}
\toprule
\textbf{Operator} & \textbf{Arity} & \textbf{Note} \\
\midrule
$+$ & 2 & Addition \\
$-$ & 2 & Subtraction \\
$*$ & 2 & Multiplication \\
$/$ & 2 & Protected division: $a\,/\,(|b|+0.0001)$ \\
\bottomrule
\end{tabular}
\end{center}

Protected division ensures the function set is closed over the reals, preventing
infinite or \texttt{NaN} outputs from divide-by-zero.

% ──────────────────────────────────────────────────────────────────────────────
\section{Initial Population Generation}

The initial population is generated using the \textbf{Grow} method:

\begin{enumerate}
  \item The root of every tree is forced to be a function node.
  \item At each non-root node, a function node is chosen with 50\% probability
        and a terminal node with 50\%, unless the maximum depth (6) is reached,
        in which case a terminal is forced.
  \item Up to three generation attempts are made per population slot; duplicate
        infix expressions are discarded and a new tree is grown to promote initial
        diversity.
\end{enumerate}

% ──────────────────────────────────────────────────────────────────────────────
\section{Fitness Function and Fitness Evaluation}

\subsection{Raw Fitness (Behavioural)}

For each training example $i$ with input vector $\mathbf{x}_i$ and true load $y_i$,
the tree $T$ is evaluated to produce $\hat{y}_i = T(\mathbf{x}_i)$. If any
prediction is \texttt{NaN} or $\pm\infty$, the individual receives fitness
\texttt{Double.MAX\_VALUE} and is excluded from the gene pool. Otherwise:

\[
\text{rawFitness}(T) = \underbrace{\frac{1}{n}\sum_{i=1}^n (\hat{y}_i - y_i)^2}_{\text{MSE}} \times (1 + 0.001 \cdot |T|)
\]

\subsection{Adjusted Fitness (Behavioural + Structural)}

Following raw fitness evaluation, structural fitness sharing (Section~3) is applied
to produce the adjusted fitness array used for selection and elitism. Individuals
in structurally crowded niches receive a fitness penalty proportional to their niche
count. Raw fitness is retained for tracking and reporting the best individual.

% ──────────────────────────────────────────────────────────────────────────────
\section{Selection Method}

\textbf{Tournament selection} with tournament size $k = 2$ is used to select
parents for crossover and mutation. Two candidates are drawn uniformly at random
(with replacement); the one with the lower \emph{adjusted} fitness wins. A small
tournament size reduces selection pressure, which, combined with the structural
fitness sharing penalty, allows structurally diverse but behaviourally competitive
individuals more opportunity to contribute offspring.

% ──────────────────────────────────────────────────────────────────────────────
\section{Genetic Operators}

\subsection{Subtree Crossover (Rate: 90\%)}

Exactly $\text{round}(980 \times 0.90) = 882$ offspring slots per generation are
produced by crossover. A random structural point in a copy of parent~1's tree is
selected; the subtree rooted there is replaced with a deep copy of a random subtree
from parent~2. One offspring is produced per crossover. Bloat control
(\texttt{trimToDepth}, max depth~6) is applied immediately after the swap.

\subsection{Subtree Mutation (Rate: 10\%)}

The remaining 98 non-elite slots per generation are produced by mutation. A random
node in a copy of the parent is selected as the mutation point; the depth remaining
from that node to the maximum depth bounds the replacement subtree. A freshly grown
subtree of at most that depth is generated and grafted in place of the old subtree.

\subsection{Exact Operator Rates via Fisher-Yates Shuffle}

To guarantee exactly 882 crossover and 98 mutation offspring per generation, a
slot-type array of length 980 is constructed (882 zeros for crossover, 98 ones for
mutation) and shuffled in-place using Fisher-Yates. This randomises the order in
which operators are applied without allowing stochastic per-individual rolls to
produce unintended ratios.

\subsection{Elitism}

The top 2\% of individuals by \emph{adjusted} fitness (20 individuals) are copied
unchanged into the next generation, ensuring structurally diverse high-quality
solutions are preserved.

% ──────────────────────────────────────────────────────────────────────────────
\section{Experimental Setup}

\subsection{Parameter Values}

\begin{center}
\begin{tabular}{@{}L{7cm}l@{}}
\toprule
\textbf{Parameter} & \textbf{Value} \\
\midrule
Population size & 1,000 \\
Max tree depth & 6 \\
Generations & 100 \\
Tournament size & 2 \\
Crossover rate & 0.90 \\
Mutation rate & 0.10 \\
Elitism rate & 2\% (20 individuals) \\
Window size & 7 \\
Training ratio & 0.80 (80/20 split) \\
Parsimony coefficient ($\lambda$) & 0.001 \\
ERC range & $[-5.0,\; 5.0]$ \\
ERC probability & 30\% \\
\midrule
\textbf{Similarity threshold ($\sigma$)} & \textbf{0.60} \\
\textbf{Sharing coefficient ($\alpha$)} & \textbf{0.01} \\
\textbf{Structural sample size} & \textbf{50} \\
\midrule
Data sample size & 100,000 rows \\
Training examples & 79,994 \\
Test examples & 19,999 \\
Independent runs & 10 \\
Random seeds & 101, 202, 303, 404, 505, 606, 707, 808, 909, 1010 \\
\bottomrule
\end{tabular}
\end{center}

\subsection{Technical Specifications}

\begin{center}
\begin{tabular}{@{}ll@{}}
\toprule
\textbf{Specification} & \textbf{Detail} \\
\midrule
Language & Java JDK 17 \\
External libraries & None (standard library only) \\
Machine & Apple MacBook (macOS) \\
Architecture & ARM64 / Apple Silicon \\
JVM & OpenJDK 17 \\
\bottomrule
\end{tabular}
\end{center}

% ──────────────────────────────────────────────────────────────────────────────
\section{Results}

All 10 runs completed successfully (10/10). Results are reported using the raw
training MSE (which excludes the parsimony component for clarity) as the primary
metric. The success predicate (train MSE $< 0.00020$) was met by all runs.
The best evolved expression was found in Run~6 (seed 606):

\begin{center}
\ttfamily\small
(((x1 + ((x2 + x4) + x3)) / 4.0562) - ((x3 * x4) - (x5 * x7)))
\end{center}

\subsection{Per-Run Results}

\begin{center}
\begin{tabular}{@{}C{1cm}C{1.5cm}C{3cm}C{3cm}C{2cm}@{}}
\toprule
\textbf{Run} & \textbf{Seed} & \textbf{Train MSE} & \textbf{Test MSE} & \textbf{Time (s)} \\
\midrule
1  & 101  & 0.00014780 & 0.00017864 & 662.14 \\
2  & 202  & 0.00014425 & 0.00017096 & 541.28 \\
3  & 303  & 0.00014899 & 0.00017794 & 780.12 \\
4  & 404  & 0.00014498 & 0.00017184 & 685.06 \\
5  & 505  & 0.00015430 & 0.00018448 & 724.11 \\
6  & \textbf{606}  & \textbf{0.00014120} & \textbf{0.00016914} & 626.64 \\
7  & 707  & 0.00014645 & 0.00017228 & 668.58 \\
8  & 808  & 0.00014402 & 0.00016999 & 743.62 \\
9  & 909  & 0.00014498 & 0.00017349 & 693.10 \\
10 & 1010 & 0.00014413 & 0.00017069 & 586.70 \\
\midrule
\textbf{Avg}    & & 0.00014611 & 0.00017394 & 671.14 \\
\textbf{StdDev} & & 0.00000342 & --         & --     \\
\textbf{Best}   & & 0.00014120 & 0.00016914 & --     \\
\bottomrule
\end{tabular}
\end{center}

\begin{figure}[H]
    \centering
    \includegraphics[width=0.95\textwidth, height=0.95\textheight, keepaspectratio]{results.jpeg}
    \caption{Terminal output from the 10-run SBGP execution (100,000 rows, seed set 101--1010).}
    \label{fig:terminal-results}
\end{figure}

\subsection{Per-Generation Statistics (Best Run)}

\begin{center}
\begin{tabular}{@{}C{1cm}C{2.5cm}C{2cm}C{1.5cm}C{1.5cm}C{2.5cm}@{}}
\toprule
\textbf{Gen} & \textbf{Avg Train MSE} & \textbf{Avg Size} & \textbf{Hits} & \textbf{Variety} & \textbf{Avg Struct Sim} \\
\midrule
10  & 0.02190210 & 4.74  & 43063 & 417 & 0.1042 \\
20  & 0.02228552 & 9.16  & 44293 & 694 & 0.0384 \\
30  & 0.02688157 & 13.85 & 44293 & 836 & 0.0380 \\
40  & 0.02068111 & 17.19 & 44416 & 917 & 0.0697 \\
50  & 0.02321578 & 21.37 & 44293 & 909 & 0.0579 \\
60  & 0.03647585 & 26.12 & 44551 & 954 & 0.1328 \\
70  & 0.02534679 & 26.63 & 44490 & 953 & 0.1670 \\
80  & 0.02623639 & 26.26 & 44579 & 960 & 0.1867 \\
90  & 0.03123360 & 26.29 & 44579 & 962 & 0.2091 \\
100 & 0.02209658 & 26.54 & 44579 & 951 & 0.2195 \\
\bottomrule
\end{tabular}
\end{center}

\subsection{Best Evolved Expression}

The best expression (Run~6, seed~606, train MSE = 0.00014120) is:

\[
\hat{y} = \frac{x_1 + x_2 + x_3 + x_4}{4.0562} - (x_3 x_4 - x_5 x_7)
\]

This expression is interpretable: the first term computes a weighted average of the
four most recent lag values (dividing by $\approx 4$ approximates a 4-step moving
average). The second term adds a small correction based on pairwise products of
older lags, capturing non-linear interaction effects between distant time steps.

\subsection{Discussion}

The Avg Struct Sim column quantifies structural diversity across the run.
In Run~6 (seed~606), values start at 0.1042 (generation~10), drop sharply to 0.0380
(generation~30) as the population becomes more structurally diverse, then rise gradually to
0.2195 by generation~100. Crucially, the similarity remains below 0.25 throughout
the entire run, confirming that SBGP successfully prevents structural convergence.

Population variety (unique infix strings) rises rapidly from 417 at generation~10 to
around 950--960 from generation~60 onwards, plateauing at near-maximum diversity.
This contrasts with canonical GP, where structural homogeneity tends to increase
steadily once a favourable operator--terminal pattern dominates.

The hit count stabilises at 44,579 / 79,994 training cases (55.7\%) from
generation~80, indicating the best individual's predictions converge to within
0.01~kW on more than half of all training examples.

Average tree size grows from 4.74 (generation~10) to 26.54 (generation~100),
reflecting the natural bloat dynamic in subtree crossover. The parsimony coefficient
($\lambda = 0.001$) limits but does not completely suppress this growth.

% ──────────────────────────────────────────────────────────────────────────────
\section{Runtimes}

Runtimes per run are reproduced in the per-run table (Section~10.1). The average
runtime across the 10 SBGP runs is 671~seconds (11.2~minutes) per run, compared
to Assignment~1's average of 298~seconds per run -- approximately $2.25\times$ slower.

The SBGP overhead arises from the structural fitness sharing step executed once per
generation: 1,000 individuals $\times$ 50 structural comparisons per individual
$= 50{,}000$ prefix-match traversals per generation. Each traversal is
$O(\min(\text{depth}_1, \text{depth}_2)) \approx O(6)$, giving
$O(300{,}000)$ additional operations per generation or $O(30{,}000{,}000)$ over
100 generations. This cost is secondary to the fitness evaluation step (79,994
tree evaluations per individual per generation), but it is non-trivial in practice.
A larger structural sample size would improve niche-count accuracy at the cost of
additional runtime; the chosen value of 50 represents a practical balance.

% ──────────────────────────────────────────────────────────────────────────────
\section{Comparison with Assignment 1}

\begin{center}
\begin{tabular}{@{}L{4cm}C{2.8cm}C{2.2cm}C{2.8cm}C{2.5cm}@{}}
\toprule
\textbf{System} & \textbf{Avg Train MSE} & \textbf{StdDev} & \textbf{Best Train MSE} & \textbf{Avg Test MSE} \\
\midrule
A1 -- Canonical GP & 0.00014812 & 0.00000658 & 0.00014257 & 0.00017889 \\
A2 -- SBGP         & 0.00014611 & 0.00000342 & 0.00014120 & 0.00017394 \\
\bottomrule
\end{tabular}
\end{center}

\subsection*{Discussion}

SBGP outperforms canonical GP on every reported metric. The key results are
summarised below:

\begin{itemize}
  \item \textbf{Average train MSE} improved by 1.4\% (0.00014812 $\to$ 0.00014611),
        corresponding to a reduction in RMSE from $\approx 0.01217$ kW to
        $\approx 0.01209$ kW.
  \item \textbf{Standard deviation} fell by 48\% (0.00000658 $\to$ 0.00000342),
        the most striking improvement. SBGP produces far more consistent results
        across independent runs, indicating that structural fitness sharing reduces
        the sensitivity of the search to the initial random seed.
  \item \textbf{Best train MSE} improved by 0.96\% (0.00014257 $\to$ 0.00014120).
  \item \textbf{Average test MSE} improved by 2.8\% (0.00017889 $\to$ 0.00017394),
        suggesting marginally better generalisation. The test/train ratio improved
        from $\approx 1.208\times$ in A1 to $\approx 1.190\times$ in A2.
  \item \textbf{Success rate:} Both systems achieved 10/10 successful runs.
\end{itemize}

These improvements are strongly attributable to the structural fitness sharing
mechanism: since all base GP parameters (population, depth, generations, operators,
seeds) are held constant between A1 and A2, the fitness sharing penalty is the
only differing component. By penalising structurally crowded individuals, SBGP
maintains greater structural diversity (Avg Struct Sim $< 0.25$ throughout all
generations), enabling the search to explore a broader variety of expression
templates. This diversity prevents the premature structural convergence observed
in canonical GP, where a single highly-fit pattern tends to dominate the population
once it proliferates.

The reduction in variance is particularly notable: SBGP's consistency across seeds
suggests the structural diversity mechanism creates a more robust search landscape,
less dependent on which structural patterns happen to emerge early in any given run.

The trade-off is runtime: SBGP requires approximately $2.25\times$ more computation
per run (avg 671~s vs avg 298~s for A1) due to the structural comparison cost.
For problems where consistency and marginal accuracy gains are valued over speed,
SBGP offers a worthwhile improvement over canonical GP.

\end{document}
